% !TEX program = xelatex
% Lernplatt - by Can Siebert
% Kurs 71 (Dig 42) - Tag 7 - Aufgabenheft
% Daten werden zu Realitaet: Event Logs, Discovery, Conformance

\documentclass[11pt,a4paper,oneside]{article}

\usepackage{polyglossia}
\setdefaultlanguage{german}
\usepackage[a4paper,margin=2.5cm]{geometry}

\usepackage{fontspec}
\defaultfontfeatures{Ligatures=TeX}
\IfFontExistsTF{Charter}{\setmainfont{Charter}}{%
  \setmainfont{Noto Serif}%
}
\IfFontExistsTF{Source Sans 3}{\setsansfont{Source Sans 3}}{%
  \IfFontExistsTF{Source Sans Pro}{\setsansfont{Source Sans Pro}}{\setsansfont{Liberation Sans}}%
}
\newcommand{\caveatfontfallback}{\itshape}
\IfFontExistsTF{Caveat}{\newfontfamily\caveatfont{Caveat}[Scale=1.15]}{\newcommand\caveatfont{\caveatfontfallback}}
\setmonofont{Liberation Mono}

\usepackage{microtype}
\usepackage{eurosym}
\linespread{1.15}

\usepackage{xcolor}
\definecolor{lpInk}{RGB}{20,28,38}
\definecolor{lpAccent}{RGB}{157,74,26}
\definecolor{lpAccentLight}{RGB}{246,228,212}
\definecolor{lpAmber}{RGB}{222,138,30}
\definecolor{lpAmberLight}{RGB}{255,243,217}
\definecolor{lpAmberBorder}{RGB}{196,118,18}
\definecolor{lpGray}{RGB}{96,104,114}
\definecolor{lpGrayLight}{RGB}{238,240,243}
\definecolor{lpGreen}{RGB}{34,120,72}
\definecolor{lpGreenLight}{RGB}{222,238,228}
\definecolor{lpGreenBorder}{RGB}{24,96,58}
\definecolor{lpL1}{RGB}{34,120,72}
\definecolor{lpL2}{RGB}{22,90,140}
\definecolor{lpL3}{RGB}{170,42,52}

\usepackage[most]{tcolorbox}
\tcbuselibrary{breakable,skins}

\newtcolorbox{infobox}[1][Hinweis]{%
  enhanced, breakable, colback=lpAccentLight, colframe=lpAccent,
  boxrule=0.6pt, arc=2pt, left=10pt,right=10pt,top=8pt,bottom=8pt,
  fonttitle=\sffamily\bfseries\color{white}, coltitle=white,
  title={#1}, attach boxed title to top left={xshift=8pt,yshift=-8pt},
  boxed title style={size=small, colback=lpAccent, colframe=lpAccent, boxrule=0.4pt, arc=1pt},
  top=14pt
}

\newtcolorbox{antwortbox}[1][Ihre Antwort]{%
  enhanced, breakable, colback=white, colframe=lpGray,
  boxrule=0.4pt, arc=2pt, left=10pt,right=10pt,top=8pt,bottom=8pt,
  fonttitle=\sffamily\bfseries\color{white}, coltitle=white,
  title={#1}, attach boxed title to top left={xshift=8pt,yshift=-8pt},
  boxed title style={size=small, colback=lpGray, colframe=lpGray, boxrule=0.4pt, arc=1pt},
  top=14pt
}

\usepackage{enumitem}
\setlist{itemsep=2pt, topsep=4pt, parsep=0pt}
\setlist[itemize,1]{label={\textbullet},leftmargin=1.6em}
\setlist[itemize,2]{label={\textendash},leftmargin=1.4em}
\setlist[enumerate,1]{label=\arabic*., leftmargin=1.8em}

\usepackage{tabularx}
\usepackage{array}
\usepackage{booktabs}
\usepackage{longtable}

\usepackage{fancyhdr}
\usepackage{lastpage}
\pagestyle{fancy}
\fancyhf{}
\renewcommand{\headrulewidth}{0.3pt}
\renewcommand{\footrulewidth}{0pt}
\fancyhead[L]{\footnotesize\sffamily\color{lpGray}Aufgabenheft \textperiodcentered{} Kurs 71 \textperiodcentered{} Tag 7}
\fancyhead[R]{\footnotesize\sffamily\color{lpGray}Discovery \textperiodcentered{} Conformance}
\fancyfoot[L]{\footnotesize\sffamily\color{lpGray}\textcopyright{} Can Siebert}
\fancyfoot[R]{\footnotesize\sffamily\color{lpGray}\thepage{} / \pageref{LastPage}}

\fancypagestyle{plain}{%
  \fancyhf{}%
  \renewcommand{\headrulewidth}{0pt}%
  \fancyfoot[C]{\footnotesize\sffamily\color{lpGray}\thepage{} / \pageref{LastPage}}%
}

\usepackage[explicit]{titlesec}
\titleformat{\section}[block]
  {\sffamily\Large\bfseries\color{lpAccent}}
  {\thesection.}{0.6em}{#1}
  [{\vspace{2pt}\color{lpAccent}\titlerule[0.6pt]}]
\titleformat{\subsection}[block]
  {\sffamily\large\bfseries\color{lpInk}}
  {\thesubsection}{0.6em}{#1}
\titlespacing*{\section}{0pt}{14pt}{8pt}
\titlespacing*{\subsection}{0pt}{10pt}{4pt}

\usepackage[hidelinks,unicode]{hyperref}

\newcommand{\akzent}[1]{{\caveatfont\color{lpAmber}#1}}
\newcommand{\fachbegriff}[1]{\textbf{#1}}
\newcommand{\quelle}[1]{{\footnotesize\sffamily\color{lpGray}(#1)}}

\newcommand{\niveaubadge}[2]{%
  \tcbox[on line, colback=#1, colframe=#1, coltext=white,
         boxrule=0pt, arc=2pt, left=5pt,right=5pt,top=1pt,bottom=1pt,
         nobeforeafter]{\sffamily\bfseries\footnotesize #2}%
}
\newcommand{\nivLi}{\niveaubadge{lpL1}{L1\,\textbullet\,Wissen}}
\newcommand{\nivLii}{\niveaubadge{lpL2}{L2\,\textbullet\,Anwendung}}
\newcommand{\nivLiii}{\niveaubadge{lpL3}{L3\,\textbullet\,Management}}

\newcommand{\zeit}[1]{%
  \tcbox[on line, colback=lpGrayLight, colframe=lpGrayLight, coltext=lpInk,
         boxrule=0pt, arc=2pt, left=5pt,right=5pt,top=1pt,bottom=1pt,
         nobeforeafter]{\sffamily\footnotesize\textbf{$\sim$\,#1}}%
}

\newcommand{\aufgabenkopf}[4]{%
  \par\vspace{6pt}
  \noindent{\sffamily\Large\bfseries\color{lpAccent}Aufgabe #1 \textendash{} #2}\par
  \vspace{2pt}
  \noindent{\color{lpAccent}\rule{\linewidth}{0.6pt}}\par
  \vspace{4pt}
  \noindent #3 \hfill \zeit{#4}\par
  \vspace{6pt}
}

\newcommand{\ytquery}[1]{%
  \par\vspace{4pt}
  \noindent{\footnotesize\sffamily\color{lpGray}\textbf{Lernvideo zum Thema:}\ %
    \href{https://studyflix.de/search?query=#1}{\textcolor{lpAccent}{Studyflix-Treffer}}\ ·\ %
    \href{https://www.youtube.com/results?search\_query=#1}{\textcolor{lpAccent}{YouTube-Treffer}}\\%
    \textit{\textcolor{lpGray}{Stichworte: #1}}}\par
}

\newcommand{\antwortlinien}[1]{%
  \par\vspace{6pt}
  \foreach \x in {1,...,#1}{%
    \noindent\makebox[\linewidth]{\dotfill}\\[6pt]
  }
}
\usepackage{pgffor}

\begin{document}

\begin{titlepage}
  \pagestyle{empty}
  \vspace*{1.5cm}
  \noindent{\sffamily\footnotesize\color{lpGray}LERNPLATT \textperiodcentered{} BY CAN SIEBERT}\\[2pt]
  \noindent{\color{lpAccent}\rule{\linewidth}{1.2pt}}\\[4pt]
  \noindent{\sffamily\footnotesize\color{lpGray}AUFGABENHEFT \textperiodcentered{} MODUL 3 \textperiodcentered{} TAG 7 VON 16 \textperiodcentered{} KURS 71 (Dig 42) \textperiodcentered{} 8.\,Juni 2026}

  \vspace{3cm}
  {\sffamily\fontsize{30}{34}\selectfont\bfseries\color{lpInk}Aufgabenheft\\Tag 7\par}

  \vspace{0.7cm}
  {\sffamily\large\color{lpGray}Daten werden zu Realit\"at \textendash{}\\Event Logs, Discovery und\\Conformance Checking.\par}

  \vspace{1.5cm}
  {\caveatfont\LARGE\color{lpAmber}11 Aufgaben \textperiodcentered{} L1 \textperiodcentered{} L2 \textperiodcentered{} L3}\par

  \vfill

  \noindent\begin{minipage}{0.55\linewidth}
    {\sffamily\footnotesize\color{lpGray}Niveau-Mix:}\\
    {\sffamily\small\color{lpInk}3\,\textsf{L1} \textperiodcentered{} 5\,\textsf{L2} \textperiodcentered{} 3\,\textsf{L3}}\\[10pt]
    {\sffamily\footnotesize\color{lpGray}Bearbeitung:}\\
    {\sffamily\small\color{lpInk}Selbstbearbeitung im Lernpfad}
  \end{minipage}\hfill
  \begin{minipage}{0.4\linewidth}
    \raggedleft
    {\sffamily\footnotesize\color{lpGray}Dozent:}\\
    {\sffamily\small\color{lpInk}Can Siebert}\\[10pt]
    {\sffamily\footnotesize\color{lpGray}Begleitend:}\\
    {\sffamily\small\color{lpInk}Lehrskript \& L\"osungsheft}
  \end{minipage}

  \vspace{0.5cm}
  \noindent{\color{lpAccent}\rule{\linewidth}{0.6pt}}
\end{titlepage}

\section*{So arbeiten Sie mit diesem Heft}
\thispagestyle{plain}

\noindent Dieses Aufgabenheft enth\"alt elf Aufgaben zu Tag 7 \textendash{} \emph{Daten werden zu Realit\"at}: Process Discovery (Ist-Prozess aus Logs), Conformance Checking (Ist vs.\ Soll) und Abweichungssteuerung. Die Aufgaben gliedern sich in drei Niveaus:

\begin{itemize}
  \item \nivLi{} \textendash{} \fachbegriff{Wissen.} Event-Log-Aufbau, Discovery-Grundlagen, Conformance-Logik. 5\,--\,10 Minuten.
  \item \nivLii{} \textendash{} \fachbegriff{Anwendung.} Happy Path ableiten, Abweichungen klassifizieren, Guardrails. 15\,--\,30 Minuten.
  \item \nivLiii{} \textendash{} \fachbegriff{Management.} Fallstudie \emph{MedTech Supplies}, Conformance Governance. 30\,--\,45 Minuten.
\end{itemize}

\noindent Jede Aufgabe enth\"alt:
\begin{itemize}
  \item eine Niveau-Markierung und einen Zeit-Richtwert,
  \item den Aufgabentext (oft mit Kontext oder Fallbeschreibung),
  \item einen Antwort-Freiraum f\"ur Ihre L\"osung,
  \item eine \emph{YouTube-Suchquery} f\"ur den begleitenden Lernpfad.
\end{itemize}

\noindent\textbf{Format der Klassifikation:} F\"ur jede Abweichung pr\"ufen Sie vier Kategorien: \emph{Compliance-Risiko} (Audit), \emph{Qualit\"ats-/Kostenrisiko}, \emph{Legitime Variante} (Innovation), \emph{Datenproblem} (Log-Fehler). Die Kategorie steuert die Reaktion.

\begin{infobox}[Hinweis]
Die Musterl\"osungen finden Sie im \emph{L\"osungsheft Tag 7}. Heute steht der Senior-Reflex im Vordergrund: \emph{Abweichungen sind Signale, nicht nur Fehler}. Welche Abweichung beweist Innovationskraft \textendash{} welche kostet sie? Pr\"ufen Sie bei jeder Aufgabe: Habe ich die Abweichung \emph{klassifiziert} (Compliance, Qualit\"at, Legitim, Daten) \textendash{} oder nur gez\"ahlt? Conformance-Steuerung ohne Klassifikation ist Reporting, nicht Steuerung.
\end{infobox}

\clearpage

\section*{Niveau L1 \textendash{} Wissen \& Reproduktion}
\addcontentsline{toc}{section}{L1 -- Wissen}

% LERNPLATT_HILFSVIDEOS
\section*{Hilfsvideos zu diesem Tag}
\addcontentsline{toc}{section}{Hilfsvideos zu diesem Tag}
\thispagestyle{plain}

\noindent Die folgenden YouTube-Erkl\"arvideos vermitteln die Inhalte, die Sie zur Bearbeitung der Aufgaben ben\"otigen. Klicken Sie auf den Titel, um das Video direkt zu \"offnen; die vollst\"andige URL steht in Klammern.

\begin{description}[style=nextline,leftmargin=18pt,labelindent=0pt,itemsep=8pt]
  \item[1.~Introduction to Process Mining — 360°-Überblick (Wil van der Aalst)] \href{https://youtu.be/aHbHQ6caJoU}{\textcolor{lpAccent}{https://youtu.be/aHbHQ6caJoU}}\\[2pt]
  {\footnotesize\color{lpGray}(URL: \nolinkurl{https://youtu.be/aHbHQ6caJoU})}
  \item[2.~Celonis Conformance Checker — Tutorial] \href{https://youtu.be/QBvjRcqsySw}{\textcolor{lpAccent}{https://youtu.be/QBvjRcqsySw}}\\[2pt]
  {\footnotesize\color{lpGray}(URL: \nolinkurl{https://youtu.be/QBvjRcqsySw})}
  \item[3.~Process Discovery — RWTH Lecture (van der Aalst)] \href{https://youtu.be/jPERVdZYgA8}{\textcolor{lpAccent}{https://youtu.be/jPERVdZYgA8}}\\[2pt]
  {\footnotesize\color{lpGray}(URL: \nolinkurl{https://youtu.be/jPERVdZYgA8})}
  \item[4.~Celonis in der Produktion — Process Mining im Einsatz] \href{https://youtu.be/RERGYW5o5xY}{\textcolor{lpAccent}{https://youtu.be/RERGYW5o5xY}}\\[2pt]
  {\footnotesize\color{lpGray}(URL: \nolinkurl{https://youtu.be/RERGYW5o5xY})}
\end{description}
\vspace{0.6em}
\noindent{\footnotesize\color{lpGray}\textit{Tipp:} \"Offnen Sie das Video, lassen Sie es im Hintergrund laufen und arbeiten Sie parallel an der Aufgabe.}

\clearpage

\aufgabenkopf{1}{Event Log: Pflichtfelder, Optional-Felder, Stolpersteine}{\nivLi}{8\,min}

\noindent Ein \fachbegriff{Event Log} ist die Grundlage f\"ur Process Mining. Listen Sie:

\begin{enumerate}
  \item Die \emph{drei Pflichtfelder} eines Event Logs: Case ID, Aktivit\"at, Timestamp.
  \item Vier \emph{Optional-Felder}, die typischerweise zu mehr Erkenntnissen f\"uhren (Ressource, Kanal, Kosten, Standort).
  \item Vier typische \emph{Stolpersteine} (fehlende/unscharfe Activities, doppelte Events, falsche Case IDs, Zeitstempel-Probleme).
\end{enumerate}

\medskip
\noindent F\"ur jeden Stolperstein:
\begin{itemize}
  \item Wie erkennen Sie ihn? (Symptom in der Discovery-Ausgabe.)
  \item Wie reagieren Sie? (Korrektur, Ausschluss, R\"uckfrage an den Quellsystem-Owner.)
\end{itemize}

\ytquery{Process Mining Event Log Required Fields Case ID Activity Timestamp Quality}

\antwortlinien{14}

\clearpage

\aufgabenkopf{2}{Discovery: Happy Path, Variante, Schleife}{\nivLi}{8\,min}

\noindent Definieren Sie pr\"agnant (je 1\,--\,2 S\"atze):

\begin{enumerate}
  \item \fachbegriff{Happy Path} \textendash{} der typischste / h\"aufigste Ablauf.
  \item \fachbegriff{Variante} \textendash{} ein abweichender Pfad (z.\,B.\ R\"uckfrage, Storno).
  \item \fachbegriff{Schleife (Rework)} \textendash{} eine Aktivit\"at wiederholt sich.
  \item \fachbegriff{Engpass} \textendash{} ein Schritt mit \"uberproportional langer Wartezeit.
\end{enumerate}

\medskip
\noindent Erkl\"aren Sie zus\"atzlich die drei Ergebnisarten von Discovery: \emph{Prozessgraph}, \emph{Variantentabelle}, \emph{Durchlaufzeiten je Schritt}.

\ytquery{Process Discovery Happy Path Variants Loops Bottleneck Beispiel}

\antwortlinien{10}

\clearpage

\aufgabenkopf{3}{Conformance: vier Abweichungstypen}{\nivLi}{10\,min}

\noindent Bei \fachbegriff{Conformance Checking} vergleichen wir Ist (Logs) mit Soll (Policy / BPMN). Listen Sie die vier zentralen Abweichungstypen:

\begin{enumerate}
  \item \emph{Fehlende Schritte} (im Soll, nicht im Ist).
  \item \emph{Zus\"atzliche Schritte} (im Ist, nicht im Soll).
  \item \emph{Falsche Reihenfolge}.
  \item \emph{Schleifen / Rework}.
\end{enumerate}

\medskip
\noindent F\"ur jeden Typ:
\begin{itemize}
  \item Beispiel aus einem Procure-to-Pay-Prozess.
  \item Welche Interpretation ist m\"oglich (Compliance-Versto{\ss} / Datenproblem / legitime Variante / Innovation)?
\end{itemize}

\noindent Schlusssatz (1 Satz): Warum ist die Interpretation einer Abweichung wichtiger als die Z\"ahlung?

\ytquery{Conformance Checking Process Mining Deviation Types Compliance Variant}

\antwortlinien{14}

\clearpage

\section*{Niveau L2 \textendash{} Anwendung \& Transfer}
\addcontentsline{toc}{section}{L2 -- Anwendung}

\aufgabenkopf{4}{Discovery ohne Tool}{\nivLii}{22\,min}

\begin{infobox}[Mini-Eventlog]
\textbf{Case 101:} Start \textrightarrow{} Anfrage\_pr\"ufen \textrightarrow{} Angebot\_erstellen \textrightarrow{} Kunde\_best\"atigt \textrightarrow{} Auftrag\_anlegen \textrightarrow{} Versand \textrightarrow{} Rechnung \textrightarrow{} Ende.

\smallskip
\textbf{Case 102:} Start \textrightarrow{} Anfrage\_pr\"ufen \textrightarrow{} R\"uckfrage \textrightarrow{} Anfrage\_pr\"ufen \textrightarrow{} Angebot\_erstellen \textrightarrow{} Kunde\_best\"atigt \textrightarrow{} Auftrag\_anlegen \textrightarrow{} Rechnung \textrightarrow{} Versand \textrightarrow{} Ende.

\smallskip
\textbf{Case 103:} Start \textrightarrow{} Anfrage\_pr\"ufen \textrightarrow{} Angebot\_erstellen \textrightarrow{} Kunde\_best\"atigt \textrightarrow{} Auftrag\_anlegen \textrightarrow{} Versand \textrightarrow{} Versand \textrightarrow{} Rechnung \textrightarrow{} Ende.

\smallskip
\textbf{Case 104:} Start \textrightarrow{} Anfrage\_pr\"ufen \textrightarrow{} Angebot\_erstellen \textrightarrow{} Kunde\_best\"atigt \textrightarrow{} Auftrag\_anlegen \textrightarrow{} Storno \textrightarrow{} Ende.
\end{infobox}

\noindent\textbf{Arbeitsauftrag:}

\begin{enumerate}
  \item Leiten Sie den \fachbegriff{Happy Path} (typischster Ablauf) ab.
  \item Identifizieren Sie \emph{drei Varianten} und beschreiben Sie je: Ursache (Hypothese), Risiko, m\"oglicher Nutzen.
  \item Markieren Sie \emph{zwei mögliche Engp\"asse} (rein logisch, ohne Zeiten) und begr\"unden Sie.
  \item Formulieren Sie \emph{drei Fragen} an die Fachseite zum Daten-/Prozessverst\"andnis.
\end{enumerate}

\ytquery{Process Mining Discovery Cases Happy Path Variants ohne Tool Beispiel}

\antwortlinien{18}

\clearpage

\aufgabenkopf{5}{Conformance ohne Tool}{\nivLii}{22\,min}

\begin{infobox}[Soll-Regeln (Policy)]
\begin{enumerate}
  \item Rechnung darf erst \emph{nach} Versand erstellt werden.
  \item R\"uckfragen m\"ussen \emph{vor} Angebotserstellung abgeschlossen sein.
  \item Doppelversand ist unzul\"assig.
  \item Storno nur \emph{vor} Versand.
\end{enumerate}

\smallskip
\textbf{Ist-Beispiele (aus Logs, siehe Aufgabe 4):}
\begin{itemize}
  \item Case 102: Rechnung \emph{vor} Versand.
  \item Case 103: Versand \emph{doppelt}.
  \item Case 102: R\"uckfrage-Schleife (ok) \textendash{} ggf.\ Angebot zu fr\"uh erstellt.
  \item Case 104: Storno (ok, vor Versand).
\end{itemize}
\end{infobox}

\noindent\textbf{Arbeitsauftrag:}

\begin{enumerate}
  \item Klassifizieren Sie jede Abweichung in: \emph{Compliance-Risiko}, \emph{Qualit\"ats-/Kostenrisiko}, \emph{Legitime Variante}, \emph{Datenproblem}.
  \item Priorisieren Sie \emph{drei Abweichungen} nach \emph{Schaden $\times$ H\"aufigkeit} (Sch\"atzung reicht).
  \item Definieren Sie zu jeder Top-Abweichung eine \emph{Gegenma{\ss}nahme} (Prozessregel, Training, Systemkontrolle, Monitoring).
  \item \textbf{Entscheidung unter Unsicherheit:} W\"ahlen Sie \emph{eine Abweichung}, die Sie bewusst \emph{zulassen} \textendash{} mit \emph{Guardrail} (Bedingung / Schwelle).
\end{enumerate}

\ytquery{Conformance Checking Deviation Classification Priorization Guardrail Mining}

\antwortlinien{18}

\clearpage

\aufgabenkopf{6}{Hypothesen vs.\ Wahrheit \textendash{} Validierung mit Fachseite}{\nivLii}{18\,min}

\begin{infobox}[Vorgabe]
Discovery liefert \emph{Hypothesen}, keine Wahrheit. Die Fachseite muss validieren \textendash{} sonst basiert die Steuerung auf Unsinn.
\end{infobox}

\noindent\textbf{Arbeitsauftrag:}

\begin{enumerate}
  \item Formulieren Sie f\"unf \emph{Validierungsfragen}, die Sie der Fachseite zur Mini-Discovery (Aufgabe 4) stellen w\"urden \textendash{} \emph{nicht} ``Stimmt das?'', sondern offene Fragen, die Annahmen pr\"ufen.
  \item Definieren Sie ein \emph{Validierungsformat}: Wer, wie lange, mit welchen Artefakten?
  \item \textbf{Anti-Muster}: Was ist die typische Falle bei Discovery-Pr\"asentationen \textendash{} und wie verhindern Sie sie?
  \item \textbf{Faustregel} (1 Satz): Wann ist eine Discovery-Aussage ``robust genug'' f\"ur eine Managemententscheidung?
\end{enumerate}

\ytquery{Process Mining Validation Stakeholder Discovery Workshop Hypothesis}

\antwortlinien{14}

\clearpage

\aufgabenkopf{7}{Bewusst zugelassene Variante \textendash{} Guardrails entwerfen}{\nivLii}{20\,min}

\begin{infobox}[Szenario]
Im \emph{MedTech Supplies}-P2P-Prozess kommen 30\,\% Rechnungen \emph{vor} Wareneingang an (Dropshipping, Serviceleistungen, fehlende Buchung). Eine ``Zero-Tolerance''-Regel w\"are operativ teuer.
\end{infobox}

\noindent\textbf{Arbeitsauftrag:}

\begin{enumerate}
  \item Formulieren Sie einen \emph{Ausnahmeprozess}: \emph{Welche Bedingung} (Schwellenwert, Lieferanten\-typ, Vertrag) macht ``Rechnung vor Wareneingang'' tempor\"ar zul\"assig?
  \item Definieren Sie \emph{drei Guardrails}: (a) Sichtbarkeit (Pflichtkennzeichen ``EXCEPTION''), (b) SLA f\"ur Nachbuchung, (c) Monitoring + Owner.
  \item \textbf{Audit-Strategie:} Wie wei{\ss} ein Auditor in 6 Monaten, dass die Ausnahme \emph{bewusst} und \emph{kontrolliert} ist \textendash{} nicht ``\"ubersehen''?
  \item Definieren Sie das \emph{Eskalations-Trigger}: ab welcher Quote eskaliert die Ausnahme zum Prozess-\"Anderungs-Antrag?
\end{enumerate}

\ytquery{Process Compliance Tolerated Exception Guardrail Audit Trail Procurement}

\antwortlinien{16}

\clearpage

\aufgabenkopf{8}{Goodhart vs.\ Echte Steuerung \textendash{} Conformance-KPI-Set}{\nivLii}{18\,min}

\begin{infobox}[Vorgabe]
F\"ur Conformance gibt es typische KPIs: \emph{Fitness-Score}, \emph{Abweichungsquote}, \emph{Top-3 Abweichungstypen}, \emph{Trend}. Jede dieser KPIs kann gespielt werden.
\end{infobox}

\noindent\textbf{Arbeitsauftrag:}

\begin{enumerate}
  \item Listen Sie f\"ur jede der vier KPIs ein typisches Manipulations-Muster (Goodhart-Falle).
  \item Definieren Sie pro KPI eine \emph{Gegenkennzahl} oder eine Kopplungs-Regel.
  \item \textbf{Anti-Reporting-Regel}: Welche KPI w\"urden Sie auf Executive-Ebene \emph{nicht} zeigen \textendash{} obwohl sie dort verlangt wird \textendash{} und wie kommunizieren Sie das?
  \item \textbf{Faustregel}: Woran erkennen Sie, dass Conformance ``echt steuert'' vs.\ ``nur Reports'' produziert?
\end{enumerate}

\ytquery{Process Mining KPI Fitness Score Conformance Anti Gaming Manipulation}

\antwortlinien{14}

\clearpage

\section*{Niveau L3 \textendash{} Management \& Entscheidungsarchitektur}
\addcontentsline{toc}{section}{L3 -- Management}

\aufgabenkopf{9}{Fallstudie \emph{MedTech Supplies} \textendash{} Procure-to-Pay Conformance}{\nivLiii}{45\,min}

\begin{infobox}[Fallstudie \emph{MedTech Supplies}]
\textbf{Unternehmen:} \emph{MedTech Supplies} (regulierte Branche, auditrelevant).\\
\textbf{Problem:} Auditoren bem\"angeln unklare Nachweise; Einkauf klagt \"uber Bürokratie. Hinweise auf \emph{Maverick Buying} (Eink\"aufe au{\ss}erhalb des Prozesses).

\smallskip
\textbf{Restriktionen:} \textbf{6 Wochen} bis Audit; Budget \textbf{70\,k\,\euro}; Logs unvollst\"andig (nur ERP + E-Mail-Tickets); Zielkonflikt Speed vs.\ Compliance.

\smallskip
\textbf{Ist-Hinweise (qualitativ):}
\begin{itemize}
  \item 30\,\% Rechnungen \emph{vor} Wareneingang.
  \item 18\,\% mit fehlender Kostenstelle.
  \item 12\,\% \"uber ``Sonderfreigabe''.
  \item 8\,\% Dubletten.
\end{itemize}
\end{infobox}

\noindent\textbf{Arbeitsauftrag:}

\begin{enumerate}
  \item \textbf{Soll-Modell P2P} (grob, 8\,--\,10 Schritte): Bedarf \textrightarrow{} Bestellung \textrightarrow{} Wareneingang \textrightarrow{} Rechnung \textrightarrow{} Freigabe \textrightarrow{} Zahlung.
  \item \textbf{Abweichungslandkarte}: pro Abweichung \emph{Risiko}, \emph{Ursache (Hypothese)}, \emph{Gegenma{\ss}nahme}, \emph{Owner}.
  \item Priorisieren Sie \emph{3 Ma{\ss}nahmen f\"ur 6 Wochen} (Auditf\"ahig) + \emph{2 Ma{\ss}nahmen f\"ur 3 Monate} (nachhaltig).
  \item \textbf{Entscheidung unter Unsicherheit:} Welche Variante (z.\,B.\ Rechnung vor Wareneingang) wird tempor\"ar zugelassen \textendash{} unter welchen Bedingungen (Guardrails)?
\end{enumerate}

\ytquery{Procure to Pay Conformance Audit Maverick Buying 3 Way Match Mining}

\antwortlinien{34}

\clearpage

\aufgabenkopf{10}{Datenverantwortung \textendash{} Logqualit\"at sichern}{\nivLiii}{25\,min}

\noindent\textbf{Diskussionsaufgabe.}

\noindent Process Mining ist nur so gut wie der Event Log. Bei \emph{MedTech} sind die Logs unvollst\"andig (nur ERP + E-Mail). Wer tr\"agt die Verantwortung f\"ur Datenqualit\"at \textendash{} und wie wird sie operationalisiert?

\medskip
\noindent\textbf{Arbeitsauftrag:}

\begin{enumerate}
  \item Definieren Sie die Rolle \fachbegriff{Data Steward} f\"ur P2P-Logs: Aufgaben (Pflege, Pr\"ufung, Reporting), Mandat (was darf sie entscheiden?), Eskalations\-pfad.
  \item Erarbeiten Sie ein \emph{Datenqualit\"ats-SLA}: 4 Kennzahlen (Vollst\"andigkeit, Aktualit\"at, Konsistenz, Eindeutigkeit) inkl.\ Ziel\-werten.
  \item \textbf{Trade-off}: Wie balancieren Sie ``perfekte Logs'' (Aufwand) gegen ``brauchbare Logs'' (Quick Insight)? Geben Sie eine Faustregel.
  \item \textbf{Audit-Beweisf\"uhrung:} Welche Artefakte zeigen einem Auditor, dass Sie Datenverantwortung \emph{ernst nehmen}?
\end{enumerate}

\ytquery{Data Steward Process Mining Event Log Quality Compliance Governance Pharma}

\antwortlinien{20}

\clearpage

\aufgabenkopf{11}{Transferprojekt \textendash{} Conformance Governance \& Decision Architecture}{\nivLiii}{45\,min}

\begin{infobox}[Rolle und Ausgangslage]
\textbf{Ihre Rolle:} Chief Compliance Officer / Chief Process Officer / Head of Audit Readiness.

\smallskip
\textbf{Auftrag:} Unternehmensweites Conformance-Programm aufsetzen, das Audit-F\"ahigkeit erh\"oht \emph{und} operative Performance nicht zerst\"ort.

\smallskip
\textbf{Restriktionen:} uneinheitliche Systeme, schwankende Log-Qualit\"at, Budget \textbf{200\,k\,\euro}, starke Standortautonomie, Zielkonflikt Compliance vs.\ Speed.
\end{infobox}

\noindent\textbf{Arbeitsauftrag:}

\begin{enumerate}
  \item \textbf{Regelwerk-Architektur}: Muss / Kann / Ausnahmeprozesse; Verantwortliche; Review-Zyklus; Dokumentation (Audit Trail).
  \item \textbf{Daten- \& Log-Standard}: Mindestfelder, Case-ID-Strategie, Datenqualit\"ats-Checks, Ownership.
  \item \textbf{Conformance Operating Model}: Rollen (Process Owner, Data Steward, Compliance Reviewer); Eskalation; Kommunikationsplan.
  \item \textbf{Portfolio der Abweichungen}: Klassifikation nach Risiko; Schwellenwerte / Guardrails; Ma{\ss}nahmen\-katalog (Prozess, Training, Systemkontrolle).
  \item \textbf{Messkonzept}: Fitness- / Abweichungsquote als Steuerung + Anti-Gaming-Mechanismen.
  \item \textbf{Roadmap 12 Wochen}: Pilotprozess w\"ahlen, Quick Wins, Skalierungsstrategie, Enablement Level 1\,--\,3.
  \item \textbf{Zwei Entscheidungen unter Unsicherheit}:
    \begin{itemize}
      \item Welche Abweichungen werden ab sofort ``\fachbegriff{Zero Tolerance}''? Begr\"undung + Risiko + Akzeptanzstrategie.
      \item Welche Datenl\"ucken akzeptieren Sie tempor\"ar \textendash{} und mit welchen Kontrollen kompensieren Sie? Guardrails + Fr\"uhwarnsignale.
    \end{itemize}
\end{enumerate}

\noindent\textbf{Bewertungskriterien:} Audit-Robustheit \textperiodcentered{} operative Tragf\"ahigkeit \textperiodcentered{} Klarheit der Verantwortung \textperiodcentered{} Entscheidungslogik.

\ytquery{Conformance Governance Decision Architecture Process Mining CCO Audit Readiness}

\antwortlinien{38}

\clearpage

\section*{Abschluss}
\thispagestyle{plain}

\noindent Sie haben alle elf Aufgaben bearbeitet. Vergleichen Sie nun Ihre Antworten mit dem \emph{L\"osungsheft Tag 7}. Pr\"ufen Sie:

\begin{itemize}
  \item Habe ich Abweichungen \emph{klassifiziert} \textendash{} oder nur gez\"ahlt?
  \item Habe ich eine bewusst zugelassene Variante \emph{mit Guardrails} ausgestattet \textendash{} oder einfach toleriert?
  \item Hat mein Conformance-KPI-Set einen wirksamen \emph{Goodhart-Schutz}?
  \item Habe ich Data Steward / Log-Verantwortung \emph{operationalisiert} \textendash{} oder nur erw\"ahnt?
\end{itemize}

\medskip
\begin{infobox}[Was Sie jetzt k\"onnen sollten]
\begin{itemize}
  \item Event Logs auf Pflichtfelder, Optional-Felder und Stolpersteine pr\"ufen.
  \item Happy Path, Varianten und Schleifen aus Mini-Logs ableiten \textendash{} ohne Tool.
  \item Conformance-Abweichungen in vier Typen klassifizieren \textendash{} und nach Schaden $\times$ H\"aufigkeit priorisieren.
  \item Eine bewusst zugelassene Variante mit Guardrails verantworten.
  \item Ein Conformance Operating Model entwerfen, das Audit-F\"ahigkeit und Performance gleichzeitig steuert.
\end{itemize}
\end{infobox}

\vspace{1cm}
\noindent{\color{lpAccent}\rule{\linewidth}{0.6pt}}\\[4pt]
{\footnotesize\sffamily\color{lpGray}Lernplatt \textperiodcentered{} by Can Siebert \textperiodcentered{} Kurs 71 (Dig 42) \textperiodcentered{} Tag 7 \textperiodcentered{} Aufgabenheft \textperiodcentered{} \textcopyright{} 2026}

\end{document}
